% Generated draft: requires mathematical and source review.
\documentclass[11pt]{article}
\usepackage[margin=25mm]{geometry}
\usepackage{fontspec}
\setmainfont{Latin Modern Roman}
\usepackage{amsmath,amssymb,mathtools,unicode-math}
\setmathfont{Latin Modern Math}
\usepackage{enumitem,needspace,xurl,hyperref}
\setlength{\parindent}{0pt}
\setlength{\parskip}{5pt}
\setlength{\emergencystretch}{3em}
\newcommand{\N}{\mathbb{N}}
\newcommand{\Z}{\mathbb{Z}}
\newcommand{\Q}{\mathbb{Q}}
\newcommand{\R}{\mathbb{R}}
\newcommand{\C}{\mathbb{C}}
\newcommand{\F}{\mathbb{F}}
\newcommand{\A}{\mathbb{A}}
\newcommand{\PP}{\mathbb{P}}
\newcommand{\E}{\mathbb{E}}
\newcommand{\eps}{\varepsilon}
\newcommand{\dd}{\,\mathrm{d}}
\begin{document}
\section*{UnsolvedMath ID: 30001416}
\textbf{Source problem code:} catalogue entry
\subsection{Definitions and mathematical statement}
For an open set \(U\subset\R^n\) and \(1\le s\le\infty\), let \(W^{1,s}(U)\) denote the space of functions in \(L^s(U)\) whose first distributional derivatives belong to \(L^s(U)\), equipped with its usual Sobolev norm. Under the bounded linear extension convention, \(U\) is a \(W^{1,s}\)-extension domain if there is a bounded linear map
\[
E_s:W^{1,s}(U)\longrightarrow W^{1,s}(\R^n)
\]
such that \((E_su)|_U=u\) almost everywhere for every \(u\in W^{1,s}(U)\). The supplied record does not explicitly specify the extension-operator convention; the following formulation uses this standard convention.

\textbf{Question 1: interpolation of extension domains.} Let \(\Omega\subset\R^n\) be a domain and let \(1\le p<q\le\infty\). If \(\Omega\) is both a \(W^{1,p}\)-extension domain and a \(W^{1,q}\)-extension domain, must it be a \(W^{1,r}\)-extension domain for every real \(r\) satisfying \(p<r<q\)? Equivalently, under the convention above, does
\[
\bigl(\exists E_p\bigr)\ \land\ \bigl(\exists E_q\bigr)
\quad\Longrightarrow\quad
\forall r\in(p,q)\;\exists E_r
\]
hold, with each operator having the corresponding extension property? No common operator at the two endpoint exponents is assumed. No boundedness or boundary regularity hypothesis on \(\Omega\) is supplied.

\textbf{Question 2: signs of a Sobolev Jacobian.} A homeomorphism is a continuous bijection with continuous inverse. For a map \(f\) with an almost everywhere defined weak derivative matrix \(Df\), write \(J_f=\det Df\). Writing \(\mathcal L^n\) for Lebesgue measure, the supplied question asks whether a Sobolev homeomorphism \(f:\Omega_1\to\Omega_2\) can satisfy both
\[
\mathcal L^n\bigl(\{x\in\Omega_1:J_f(x)>0\}\bigr)>0,
\qquad
\mathcal L^n\bigl(\{x\in\Omega_1:J_f(x)<0\}\bigr)>0.
\]
Here the intended setting appears to be domains \(\Omega_1,\Omega_2\subset\R^n\), but the record does not specify the dimension range, the Sobolev integrability exponent, or whether Sobolev membership is global or local. For example, membership in \(W^{1,s}_{\mathrm{loc}}(\Omega_1;\R^n)\), meaning Sobolev membership on every relatively compact open subset, would provide a weak derivative matrix for \(s\ge1\); this is an explanatory framework, not a supplied choice of hypotheses. The sign conditions above state the requested phenomenon precisely, but the class of admissible maps remains unresolved. No orientation assumption or regularity of the inverse beyond continuity is specified.
\subsection{Short English statement}
The first question asks whether extension of first-order Sobolev functions at two integrability exponents forces extension at every intermediate exponent. The second asks whether the Jacobian of a Sobolev homeomorphism can be positive and negative on sets of positive measure. Its dimension and Sobolev regularity requirements are missing, so a fully quantified formulation requires clarification from the source. The supplied record labels the combined problem partially solved; this is supplied curation, not an independently verified status.
\subsection{Sources}
\begin{itemize}
\item Supplied local UnsolvedMath registry record, permanent ID \(30001416\), title \emph{Interpolation of Sobolev Extension Domains and Jacobian Signs}, source problem code \texttt{catalogue entry}, in the Oberwolfach Reports Open Problems collection. The record gives proposed year 2010 and names no individual proposer.
\item \emph{Optimal Constants in the Theory of Sobolev Spaces and PDEs} (2010), Oberwolfach Reports, pp.~325--371, DOI \texttt{10.4171/owr/2010/08}. Supplied source report and literature reference: \url{https://doi.org/10.4171/owr/2010/08}. The supplied background describes Oberwolfach Reports as documenting mathematical workshops and their open-problem discussions; no precise problem page is supplied.
\item Goldstein--Hajlasz, 2018, supplied literature reference: \url{https://arxiv.org/abs/1812.11888}. No paper title or theorem number is supplied.
\item Literature assessment dated \texttt{2026-08-21}: the supplied curation labels the record partially solved and states that Goldstein and Hajlasz reduce the unresolved Sobolev-homeomorphism case to the critical exponent and prove nonnegativity under additional hypotheses, while the extension-domain interpolation question remains unresolved in the stated generality. This is supplied dated curation, not independently verified. The critical exponent and additional hypotheses are not specified in the record.
\item The registry labels the statement \texttt{corrected\_verified} and reports a manual semantic, typographic, and LaTeX review of its cleanup. These are supplied editorial assertions and do not establish mathematical completeness or independent source verification.
\end{itemize}
\end{document}
